\documentclass[aspectratio=169,10pt]{beamer}
\usepackage{beamer_theme_cienciadados2026}

\title[Aula 02 - Linguagens e Paradigmas]{Ciência dos Dados I: Análise Exploratória}
\subtitle{Aula 02: Linguagens e Paradigmas de Programação}
\author{Prof. Dr. Ney Lemke}
\institute{UNESP -- Instituto de Biociências de Botucatu}
\date{Versão 2026}

% Rodapé com margem inferior ampliada para projeção.
\setbeamertemplate{footline}{%
  \begin{beamercolorbox}[wd=\paperwidth,ht=3.2ex,dp=1.5ex,leftskip=1cm,rightskip=1cm]{palette secondary}
    \footnotesize
    \textbf{Ciência dos Dados (2026)} \hfill Aula 02 - Linguagens e Paradigmas \hfill \insertframenumber{} / \inserttotalframenumber
  \end{beamercolorbox}%
}

\begin{document}

\frame{\titlepage}

\begin{frame}{Percurso da aula}
  \tableofcontents
\end{frame}

\section{Linguagens de programação}

\begin{frame}{Linguagem, paradigma e ferramenta}
  \begin{columns}[T,onlytextwidth]
    \begin{column}{0.48\textwidth}
      \begin{block}{Linguagem}
        Sistema de regras para expressar computações: sintaxe, semântica e formas de representar valores e operações.
      \end{block}
      \begin{block}{Paradigma}
        Uma maneira de organizar o raciocínio e estruturar a solução de um problema.
      \end{block}
    \end{column}
    \begin{column}{0.48\textwidth}
      \begin{exampleblock}{Implementação}
        Compilador, interpretador, máquina virtual e JIT participam da tradução e execução do programa.
      \end{exampleblock}
      \begin{alertblock}{Ferramenta}
        Jupyter, editor, biblioteca e ambiente de desenvolvimento auxiliam a escrever, executar ou analisar programas; não são paradigmas.
      \end{alertblock}
    \end{column}
  \end{columns}
  \vspace{0.4em}
  \centering\textbf{Exemplo: Python é uma linguagem; Jupyter é um ambiente de trabalho.}
\end{frame}

\begin{frame}{Um panorama de linguagens}
  \small
  \begin{tabular}{p{0.32\textwidth}p{0.58\textwidth}}
    \toprule
    \textbf{Exemplos} & \textbf{Contextos frequentes (não exclusivos)} \\
    \midrule
    C, Fortran & Sistemas e aplicações científicas de alto desempenho. \\
    \addlinespace[0.15em]
    Lisp & Programação simbólica, linguagens funcionais e ambientes interativos. \\
    \addlinespace[0.15em]
    Python, Perl, Ruby, Lua & Automação, scripts e aplicações de propósito geral. \\
    \addlinespace[0.15em]
    R, MATLAB, \mbox{Mathematica}, Julia & Estatística, matemática, engenharia e computação científica. \\
    \addlinespace[0.15em]
    \shortstack[l]{Java, JavaScript,\\ Scala, Clojure, Go} & Aplicações, serviços, web e sistemas de diferentes escalas. \\
    \bottomrule
  \end{tabular}
  \vspace{0.7em}
  \begin{alertblock}{Como ler esta tabela}
    A associação indica usos comuns, não uma restrição. Linguagens evoluem, compartilham recursos e podem atender a vários domínios.
  \end{alertblock}
\end{frame}

\begin{frame}{C e Fortran: controle e computação científica}
  \begin{columns}[T,onlytextwidth]
    \begin{column}{0.48\textwidth}
      \begin{block}{C}
        \begin{itemize}
          \item Linguagem de propósito geral com controle explícito de recursos.
          \item Presente em sistemas, bibliotecas e código de alto desempenho.
          \item Exige atenção a detalhes de implementação e gerenciamento de memória.
        \end{itemize}
      \end{block}
    \end{column}
    \begin{column}{0.48\textwidth}
      \begin{exampleblock}{Fortran}
        \begin{itemize}
          \item Desenvolvida com forte relação com o cálculo numérico e científico.
          \item Ecossistema maduro para aplicações científicas e computação de alto desempenho.
          \item A linguagem evoluiu; não se limita ao estilo de seus primeiros programas.
        \end{itemize}
      \end{exampleblock}
    \end{column}
  \end{columns}
  \vspace{0.3em}
  \centering\small Desempenho depende também de algoritmos, compiladores, bibliotecas e hardware — não apenas da linguagem.
\end{frame}

\begin{frame}{Lisp e linguagens científicas}
  \begin{columns}[T,onlytextwidth]
    \begin{column}{0.48\textwidth}
      \begin{block}{Lisp}
        \begin{itemize}
          \item Associada à computação simbólica e à programação funcional.
          \item O ciclo interativo REPL (ler, avaliar e imprimir resultados) favorece experimentação.
          \item Estruturas de dados e código podem ser manipulados de formas expressivas.
        \end{itemize}
      \end{block}
    \end{column}
    \begin{column}{0.48\textwidth}
      \begin{exampleblock}{R, MATLAB, Mathematica e Julia}
        \begin{itemize}
          \item Ambientes e ecossistemas voltados a tarefas numéricas ou científicas.
          \item R é amplamente usado em estatística; MATLAB em computação numérica e engenharia; Mathematica em matemática simbólica; Julia em computação científica.
          \item A escolha depende de bibliotecas, licença, equipe e requisitos do projeto.
        \end{itemize}
      \end{exampleblock}
    \end{column}
  \end{columns}
\end{frame}

\begin{frame}{Java, JavaScript, Python e outras linguagens}
  \begin{itemize}
    \item \textbf{Java}: linguagem de propósito geral, fortemente associada ao desenvolvimento baseado em classes e à máquina virtual Java.
    \item \textbf{JavaScript}: linguagem distinta de Java, usada em navegadores e também em outros ambientes de execução.
    \item \textbf{Python}: linguagem de propósito geral, com ecossistema amplo e suporte a vários paradigmas; muito usada em análise de dados.
    \item \textbf{Perl, Ruby, Lua, Scala, Clojure e Go}: exemplos de linguagens com diferentes escolhas de sintaxe, execução, tipos e ecossistemas.
  \end{itemize}
  \begin{alertblock}{Evite rótulos simplificadores}
    “Linguagem de script”, “compilada” ou “orientada a objetos” não resume por completo uma linguagem moderna nem todos os seus usos.
  \end{alertblock}
\end{frame}

\section{Paradigmas de programação}

\begin{frame}{Programação estruturada / imperativa}
  \begin{block}{Ideia central}
    Expressar a solução como instruções e etapas explícitas, organizadas por estruturas de controle.
  \end{block}
  \begin{columns}[T,onlytextwidth]
    \begin{column}{0.32\textwidth}
      \begin{exampleblock}{Sequência}
        Executar instruções em uma ordem definida.
      \end{exampleblock}
    \end{column}
    \begin{column}{0.32\textwidth}
      \begin{exampleblock}{Seleção}
        Escolher caminhos com condicionais.
      \end{exampleblock}
    \end{column}
    \begin{column}{0.32\textwidth}
      \begin{exampleblock}{Repetição}
        Repetir operações com laços.
      \end{exampleblock}
    \end{column}
  \end{columns}
  \begin{itemize}
    \item Funções e sub-rotinas ajudam a modularizar e reduzir duplicação.
    \item Vetores, matrizes, registros e outras estruturas organizam os dados.
    \item O estado mutável pode tornar o fluxo difícil de acompanhar se não for controlado.
  \end{itemize}
\end{frame}

\begin{frame}{Orientação a objetos}
  \begin{block}{Ideia central}
    Organizar parte do sistema em objetos que reúnem estado e comportamento.
  \end{block}
  \begin{columns}[T,onlytextwidth]
    \begin{column}{0.48\textwidth}
      \begin{itemize}
        \item \textbf{Atributos}: dados associados a um objeto.
        \item \textbf{Métodos}: operações relacionadas ao objeto.
        \item \textbf{Classes}: estruturas para definir objetos e seus comportamentos.
      \end{itemize}
    \end{column}
    \begin{column}{0.48\textwidth}
      \begin{exampleblock}{Quando pode ajudar?}
        Sistemas com entidades, responsabilidades e interações que precisam ser organizadas e mantidas ao longo do tempo.
      \end{exampleblock}
      \begin{alertblock}{Cuidado}
        Classes e hierarquias podem acrescentar complexidade sem benefício; herança não é a única forma de reutilização.
      \end{alertblock}
    \end{column}
  \end{columns}
\end{frame}

\begin{frame}{Programação funcional}
  \begin{block}{Ideia central}
    Expressar computações como composição de funções e transformações de valores.
  \end{block}
  \begin{itemize}
    \item \textbf{Função pura}: para a mesma entrada, produz o mesmo resultado e não altera estado externo.
    \item \textbf{Imutabilidade}: em vez de modificar um dado compartilhado, produz-se um novo valor.
    \item \textbf{Composição}: funções menores são conectadas para formar uma transformação maior.
    \item \texttt{map} transforma cada elemento; operações de redução agregam vários valores em um resultado.
  \end{itemize}
  \begin{exampleblock}{Onde aparece?}
    Haskell e Lisp são exemplos associados ao estilo funcional; Python, R e outras linguagens também oferecem recursos funcionais.
  \end{exampleblock}
\end{frame}

\begin{frame}{Declarativo, dirigido por eventos e reativo}
  \begin{columns}[T,onlytextwidth]
    \begin{column}{0.48\textwidth}
      \begin{block}{Declarativo}
        Descrever o resultado, as relações ou as restrições, sem definir todo o algoritmo.\par\small Exemplos: SQL e Prolog.
      \end{block}
      \begin{exampleblock}{Vantagem e limite}
        Abstrai detalhes, mas parte do controle de execução fica a cargo do sistema ou do domínio.
      \end{exampleblock}
    \end{column}
    \begin{column}{0.48\textwidth}
      \begin{block}{Dirigido por eventos}
        Responde a eventos, como ações de usuário, mensagens ou mudanças de estado.
      \end{block}
      \begin{exampleblock}{Reativo}
        Propaga mudanças e responde a fluxos de dados ou eventos ao longo do tempo.
      \end{exampleblock}
    \end{column}
  \end{columns}
  \vspace{0.5em}
  \centering\small Esses estilos podem coexistir com abordagens imperativas, funcionais e orientadas a objetos.
\end{frame}

\begin{frame}[fragile]{Uma tarefa, três formas de expressar a solução}
  \small Tarefa: selecionar países com mais de 200 registros e obter seus nomes.
  \begin{columns}[T,onlytextwidth]
    \begin{column}{0.49\textwidth}
      \begin{block}{Estruturada / imperativa}
\begin{lstlisting}
selecionar = lista_vazia
para registro em tabela:
    se registro.quantidade > 200:
        adicionar registro.pais
\end{lstlisting}
      \end{block}
    \end{column}
    \begin{column}{0.49\textwidth}
      \begin{exampleblock}{Funcional}
\begin{lstlisting}
tabela
  -> filtrar quantidade > 200
  -> extrair pais
\end{lstlisting}
      \end{exampleblock}
    \end{column}
  \end{columns}
  \begin{block}{Declarativa (SQL)}
\begin{lstlisting}
SELECT pais FROM catalogo WHERE quantidade > 200;
\end{lstlisting}
  \end{block}
  \small Perguntas: qual expressa os passos? Qual descreve diretamente a condição? O que precisa ser igual para todas as respostas?
\end{frame}

\begin{frame}{Python: uma linguagem multiparadigma}
  \begin{columns}[T,onlytextwidth]
    \begin{column}{0.48\textwidth}
      \begin{block}{Pode combinar estilos}
        \begin{itemize}
          \item imperativo e estruturado;
          \item orientação a objetos;
          \item composição funcional;
          \item bibliotecas com outras abstrações.
        \end{itemize}
      \end{block}
    \end{column}
    \begin{column}{0.48\textwidth}
      \begin{exampleblock}{A escolha é contextual}
        Um mesmo projeto pode usar classes para organizar componentes, funções para transformar dados e instruções imperativas para coordenar etapas.
      \end{exampleblock}
    \end{column}
  \end{columns}
  \begin{alertblock}{Não existe “o paradigma correto” em abstrato}
    Compare legibilidade, estado, testes, manutenção, desempenho e familiaridade da equipe.
  \end{alertblock}
\end{frame}

\section{Escolha e organização do desenvolvimento}

\begin{frame}{Qual linguagem é a melhor?}
  \begin{block}{A pergunta mais útil}
    Qual linguagem, paradigma e ecossistema atendem melhor a este problema, sob estas restrições?
  \end{block}
  \begin{itemize}
    \item Que bibliotecas e ferramentas a tarefa exige?
    \item Que desempenho e recursos de hardware são necessários?
    \item Quem manterá o software e qual experiência essa equipe possui?
    \item Há requisitos de interoperabilidade, licença ou implantação?
    \item O código precisa ser exploratório, reutilizável, distribuído ou de longa duração?
  \end{itemize}
  \centering\textbf{Toda escolha envolve compromissos; nenhuma linguagem é universalmente melhor.}
\end{frame}

\begin{frame}{Catedral e Bazar: outra dimensão}
  \begin{columns}[T,onlytextwidth]
    \begin{column}{0.48\textwidth}
      \begin{block}{Catedral}
        Metáfora para desenvolvimento mais centralizado e planejado, com coordenação concentrada e entregas mais controladas.
      \end{block}
    \end{column}
    \begin{column}{0.48\textwidth}
      \begin{exampleblock}{Bazar}
        Metáfora para desenvolvimento aberto e distribuído, com contribuições, discussão e revisão da comunidade.
      \end{exampleblock}
    \end{column}
  \end{columns}
  \begin{alertblock}{Não confundir}
    Catedral/Bazar descrevem modelos de organização do desenvolvimento, não paradigmas de programação. Um projeto de qualquer modelo pode combinar vários estilos de código.
  \end{alertblock}
  \footnotesize Inspirado no ensaio de Eric S. Raymond, \textit{The Cathedral and the Bazaar} (1999).
\end{frame}

\begin{frame}{Atividade: escolha, investigue e justifique}
  \begin{enumerate}
    \item Escolha um cenário: análise de dados em notebook, consulta a banco, cálculo numérico ou sistema mantido por uma equipe.
    \item Em dupla, formulem um prompt próprio para investigar linguagens, paradigmas ou ecossistemas relevantes ao cenário.
    \item Consultem as fontes no NotebookLM e verifiquem as referências no material.
    \item Registrem o prompt inicial, eventual reformulação, evidência consultada e conclusão.
    \item Justifiquem uma escolha e indiquem pelo menos um limite ou custo.
  \end{enumerate}
  \begin{exampleblock}{Entrega}
    Cenário + escolha + critérios + limitação + fonte conferida.
  \end{exampleblock}
\end{frame}

\section{Síntese}

\begin{frame}{Síntese e referências}
  \begin{itemize}
    \item Linguagem, paradigma, ferramenta e estratégia de execução são conceitos distintos.
    \item Estruturada, orientada a objetos, funcional e declarativa são abordagens que ajudam a organizar soluções.
    \item Linguagens podem combinar paradigmas; usos comuns não são fronteiras rígidas.
    \item A escolha depende de requisitos, ecossistema, equipe e manutenção.
    \item Catedral/Bazar descrevem organização do desenvolvimento, não estrutura do código.
  \end{itemize}
  \vspace{0.3em}
  \footnotesize
  \textbf{Material de base:} CienciadosDados2024, `Aula03/Paradigmas.ipynb`, `.Paradigmas.ipynb.sage-jupyter2`, `.Linguagens de Programação.ipynb.sage-jupyter2` e `Aula04/Linguagens de Programação.ipynb`.\\
  Sebesta, R. W. \textit{Concepts of Programming Languages}, 12th ed., Pearson, 2019.\\
  Python Software Foundation. \textit{Functional Programming HOWTO}.\\
  Raymond, E. S. \textit{The Cathedral and the Bazaar}, 1999.
\end{frame}

\end{document}
